%%%%%%%%%%  scan idx 81  |  folio 68  %%%%%%%%%%
Il
suffit d'utiliser (2.37) écrit sous la forme~:
\begin{equation}\tag{2.37 bis}
 \widetilde c(i_!(x))=1+i_*\left(\frac{\widetilde c(x)*\lambda_{-1}(c(N))-1}{(-1)^qN^q}\right),
\end{equation}
la formule explicite (1.29), la formule (iv bis) et enfin (1.30).

On peut écrire (2.38) sous la forme
\begin{equation}\tag{2.38 bis}
 \ch_X(F)=i_*\left(\ch_Y(F)\,\mathcal T(T_{X/Y})^{-1}\right),
\end{equation}
où $F$ est un faisceau algébrique cohérent sur $Y$ (considéré au premier membre comme un faisceau sur $X$ nul en dehors de $Y$), $\ch_X$ et $\ch_Y$ sont les homomorphismes de Chern pris respectivement sur $X$ et $Y$, et $T_{X/Y}$ est le fibré normal à $Y$ dans $X$.

\textbf{Remarques :} Contrairement à (2.38), la formule (2.37) est une «~formule sans dénominateur~». Elle n'a été démontrée (même en caractéristique $0$) qu'en se plaçant dans $G(\underline K(X))$ et non dans $A(X)$, ce qui n'est peut-être pas la même chose. Elle est certainement vraie en toute caractéristique et dans $A(X)$. Notons que l'on vérifie tout de suite que les images des deux membres par $i^*$ sont les mêmes, d'où résulte en appliquant $i_*$ que le produit de leur différence par $N^q$ est nulle. Cela suggère une méthode de démonstration complètement différente, en essayant d'obtenir $Y$ comme image réciproque transversale d'une sous-variété $Y'$ dans une variété $X'$, qui soit telle qu'en dimension $\leq n=\dim X$, la classe de $Y'$ dans $A(X')$ ne soit pas diviseur de $0$\dots. Malheureusement on ne peut espérer que $X'$ soit une grasmanienne ou quelque chose d'approchant.

On notera que, si l'on néglige les éléments de torsion, la formule (2.38) est équivalente à (2.37).




% Additional macro definitions
\providecommand{\cS}{\mathcal S}
\providecommand{\cD}{\mathcal D}
\providecommand{\cE}{\mathscr E}
\providecommand{\cG}{\mathscr G}
\providecommand{\cL}{\mathscr L}
\providecommand{\cM}{\mathscr M}
\providecommand{\cP}{\mathscr P}
\providecommand{\cK}{\mathcal K}
\providecommand{\Db}{\mathrm D}
\providecommand{\Kb}{\mathrm K}
\providecommand{\Trc}{\operatorname{Tr}}
\providecommand{\ob}{\operatorname{ob}}
\providecommand{\Coh}{\operatorname{Coh}}
\providecommand{\Cone}{\operatorname{Cone}}
\providecommand{\im}{\operatorname{Im}}
\providecommand{\Id}{\operatorname{Id}}
\providecommand{\RR}{\mathrm{R.R.}}
\providecommand{\RRR}{\mathrm{R.R.R.}}
\providecommand{\GrR}{\mathrm{G}}
\providecommand{\pr}{\operatorname{pr}}


